\documentclass[tikz,border=6pt]{standalone}
\usepackage[no-math]{fontspec}
\usepackage[brazil]{babel}
\usepackage{amsmath,amssymb}
\IfFileExists{C:/Windows/Fonts/Lato_Regular.ttf}{%
  \setmainfont{Lato}[
    Path = C:/Windows/Fonts/,
    UprightFont = Lato_Regular.ttf,
    BoldFont = Lato_Bold.ttf,
    ItalicFont = Lato_Italic.ttf,
    BoldItalicFont = Lato_BoldItalic.ttf
  ]
  \setsansfont{Lato}[
    Path = C:/Windows/Fonts/,
    UprightFont = Lato_Regular.ttf,
    BoldFont = Lato_Bold.ttf,
    ItalicFont = Lato_Italic.ttf,
    BoldItalicFont = Lato_BoldItalic.ttf
  ]
}{%
  \RequirePackage[default]{lato}
}
\renewcommand{\familydefault}{\sfdefault}

\usepackage{xcolor}
\usepackage{tikz}
\usepackage{pgfplots}
\pgfplotsset{compat=1.18}
\usetikzlibrary{arrows.meta,positioning,shapes.geometric,fit,backgrounds,calc,decorations.pathreplacing,patterns}

% Paleta institucional UEFS / CHF662
\definecolor{uefsNavy}{HTML}{1B2A4A}
\definecolor{uefsTeal}{HTML}{0E7C7B}
\definecolor{uefsForest}{HTML}{2D6A4F}
\definecolor{uefsGold}{HTML}{C4983B}
\definecolor{uefsGoldDark}{HTML}{8F6B1E}
\definecolor{uefsAlert}{HTML}{B02A37}
\definecolor{uefsOrange}{HTML}{D95F02}
\definecolor{uefsInk}{HTML}{22252A}
\definecolor{uefsMuted}{HTML}{5C6470}
\definecolor{uefsWater}{HTML}{2E86AB}
\definecolor{uefsSand}{HTML}{D9B382}
\definecolor{bgNavy}{HTML}{F2F5F9}
\definecolor{bgTeal}{HTML}{E9F5F4}
\definecolor{bgForest}{HTML}{EEF5F0}
\definecolor{bgGold}{HTML}{FCF7EB}
\definecolor{bgAlert}{HTML}{FDF2F2}
\definecolor{bgGray}{HTML}{F8FAFC}
\definecolor{lineGray}{HTML}{CFD8DC}

\tikzset{
  >=Stealth,
  aaCard/.style={draw=#1, fill=white, rounded corners=4pt, line width=1pt, inner sep=6pt, align=center, font=\small},
  aaBox/.style={draw=#1, rounded corners=4pt, line width=1.1pt, inner sep=6pt, align=center, font=\footnotesize\sffamily, text=uefsInk},
  aaTitle/.style={font=\normalsize\bfseries\sffamily, text=uefsNavy, align=center},
  aaSub/.style={font=\footnotesize\sffamily, text=uefsMuted, align=center},
  aaArrow/.style={->, >=Stealth, line width=1.4pt},
  aaDArrow/.style={<->, >=Stealth, line width=1.4pt}
}


\begin{document}
\begin{tikzpicture}[font=\sffamily]

% =========================================================
% PAINEL (a) Erro: Círculos Concêntricos Arbitrários
% =========================================================
\begin{scope}[shift={(0,0)}]
  \draw[draw=uefsOrange, fill=bgGray, rounded corners=6pt, line width=1.2pt] (0,0) rectangle (9.6, 8.2);
  \node[font=\normalsize\bfseries, text=uefsOrange!90!black] at (4.8, 7.6) {(a) Erro: Círculos Concêntricos Arbitrários};

  % Divisores reais ignorados (polígono fantasma ao fundo)
  \draw[dashed, line width=1.3pt, uefsNavy!30, fill=uefsNavy!4]
    (1.8, 6.2) -- (6.9, 5.9) -- (7.8, 2.9) -- (4.8, 1.05) -- (1.1, 1.45) -- cycle;

  % Rede de drenagem ao fundo
  \draw[line width=1.3pt, uefsWater!65] (4.5, 5.6) -- (4.3, 4.1) -- (5.05, 1.15);
  \draw[line width=1pt, uefsWater!65] (3.1, 4.4) -- (4.35, 3.85);
  \draw[line width=1pt, uefsWater!65] (5.7, 3.3) -- (4.65, 2.75);

  % Círculos concêntricos arbitrários
  \draw[dashed, line width=1.6pt, uefsOrange] (4.5, 3.9) circle (1.45cm);
  \draw[dashed, line width=1.6pt, uefsOrange] (4.5, 3.9) circle (2.35cm);

  % ADA no centro
  \node[fill=uefsOrange!90!black, text=white, font=\footnotesize\bfseries,
        minimum width=0.9cm, minimum height=0.65cm, rounded corners=2pt]
    at (4.5, 3.9) {ADA};

  \node[font=\footnotesize\bfseries, text=uefsOrange!90!black, align=center]
    at (4.8, 0.55) {Ignora divisores topográficos e drenagem a jusante};
\end{scope}

% =========================================================
% PAINEL (b) Delimitação Técnica por Divisores Topográficos
% =========================================================
\begin{scope}[shift={(10.2, 0)}]
  \draw[draw=uefsForest, fill=bgGray, rounded corners=6pt, line width=1.2pt] (0,0) rectangle (9.6, 8.2);
  \node[font=\normalsize\bfseries, text=uefsForest] at (4.8, 7.6) {(b) Delimitação Técnica por Divisores Topográficos};

  % AII (Bacia regional / municípios) - retângulo externo pontilhado amplo
  \draw[dotted, line width=1.2pt, uefsNavy, fill=uefsWater!7, rounded corners=4pt]
    (0.85, 0.95) rectangle (8.75, 6.85);

  % Rótulo da AII em espaço livre no topo esquerdo
  \node[anchor=west, font=\footnotesize\bfseries, text=uefsNavy]
    at (1.15, 6.42) {AII (Bacia regional / municípios)};

  % AID Física (Divisores topográficos) - polígono da microbacia
  \draw[dashed, line width=1.7pt, uefsTeal, fill=uefsWater!13]
    (2.35, 5.55) -- (5.8, 5.75) -- (7.45, 3.7) -- (6.35, 1.45) -- (2.15, 1.55) -- (1.45, 3.65) -- cycle;

  % Rede de drenagem dentro da bacia
  \draw[line width=1.8pt, uefsWater!90!black] (4.3, 5.05) -- (4.7, 3.45) -- (5.1, 1.48);
  \draw[line width=1.2pt, uefsWater!90!black] (2.9, 4.35) -- (4.52, 3.95);
  \draw[line width=1.2pt, uefsWater!90!black] (6.1, 3.55) -- (4.85, 2.75);

  % ADA (Área Diretamente Afetada)
  \node[fill=uefsOrange!90!black, text=white, font=\footnotesize\bfseries,
        minimum width=0.95cm, minimum height=0.65cm, rounded corners=2pt]
    (ada) at (4.35, 4.55) {ADA};

  % Rótulo AID Física posicionado em espaço totalmente livre à direita/fora da linha tracejada ou dentro sem cruzar bordas
  \node[anchor=west, align=left, font=\footnotesize\bfseries, text=uefsTeal, fill=uefsWater!7, inner sep=2pt]
    at (6.05, 5.95) {AID Física\\[-1pt]{\scriptsize(Divisores topográficos)}};

  % Pequena seta indicadora limpa do rótulo para o divisor tracejado
  \draw[->, >=Stealth, line width=0.9pt, uefsTeal] (6.02, 5.82) -- (5.65, 5.65);

  \node[font=\footnotesize\bfseries, text=uefsForest, align=center]
    at (4.8, 0.55) {Acompanha o escoamento hídrico até o exutório};
\end{scope}

\end{tikzpicture}
\end{document}
